% concatenated from article.tex
\documentclass[a4paper,11pt]{article}

\usepackage{authblk}
\usepackage[utf8]{inputenc}
\usepackage[T1]{fontenc}
\usepackage{indentfirst}
\usepackage{mathtools}
\usepackage{verbatim}
\usepackage{url}
\usepackage{booktabs}
\usepackage{amsfonts}
\usepackage{amsmath}
\usepackage{amssymb}
\usepackage{amsthm}
\usepackage[a4paper,margin=2.5cm]{geometry}
\usepackage{enumitem}
\usepackage{xcolor}
\usepackage{mathrsfs}
\usepackage{fancyhdr}
\usepackage{listings}
\usepackage[linktocpage=true]{hyperref}
\hypersetup{
    colorlinks=true,
    linkcolor=blue,
    citecolor=violet,
    urlcolor=blue
}
\usepackage{graphicx}
\usepackage[makeroom]{cancel}
\usepackage{tikz}

\sloppy

\theoremstyle{plain}
\newtheorem{theorem}{Theorem}[section]
\newtheorem{proposition}[theorem]{Proposition}
\newtheorem{corollary}[theorem]{Corollary}
\newtheorem{lemma}[theorem]{Lemma}

\theoremstyle{definition}
\newtheorem{definition}[theorem]{Definition}
\newtheorem{example}[theorem]{Example}
\newtheorem{exercise}[theorem]{Exercise}
\newtheorem{problem}[theorem]{Problem}

\theoremstyle{remark}
\newtheorem{remark}[theorem]{Remark}

\DeclareMathOperator*{\essinf}{ess\,inf}
\DeclareMathOperator*{\esssup}{ess\,sup}
\DeclareMathOperator{\Tr}{Tr}
\DeclareMathOperator{\Gram}{Gram}
\DeclareMathOperator{\Span}{span}
\DeclareMathOperator{\Dim}{dim}
\DeclareMathOperator{\Terms}{Terms}
\DeclareMathOperator{\cl}{cl}
\DeclareMathOperator{\Ker}{Ker}
\DeclareMathOperator{\Ran}{Ran}
\DeclareMathOperator{\MultiLinForm}{MultiLinForm}
\DeclareMathOperator{\Gr}{Gr}
\DeclareMathOperator{\diag}{diag}

\makeatletter
\def\keywords{\xdef\@thefnmark{}\@footnotetext}
\makeatother

\newcommand*\circled[1]{\tikz[baseline=(char.base)]{\node[shape=circle,draw,inner sep=2pt] (char) {#1};}}

\title{Frames generated by the functional calculus and function frames of a normal operator}
\author{Nizar El Idrissi}

\newcommand{\Addresses}{{% 
  \bigskip
  \footnotesize
  \textbf{Nizar El Idrissi.}
  \par\nopagebreak Laboratoire : Équations aux dérivées partielles, Algèbre et Géométrie spectrales.
  \par\nopagebreak
  Département de mathématiques, faculté des sciences, université Ibn Tofail, 14000 Kénitra.\par\nopagebreak 
  \textit{E-mail address} : \texttt{nizar.elidrissi@uit.ac.ma}
}}

\begin{document}
\maketitle


\begin{abstract}
In this article, we prove that sequences generated by the functional calculus $(f(T)(e_n))_{n \in \mathbb{N}}$ can be equivalently written as function sequences $(f_n(T) g)_{n \in \mathbb{N}}$, when $T$ is normal and $g$ a cyclic vector for $T$. Here, $(e_n)_{n \in \mathbb{N}}$ is a sequence of vectors, $T$ is a bounded normal operator, $f$ and $(f_n)_{n \in \mathbb{N}}$ are functions defined on a neighborhood of the spectrum $\sigma(T)$, and $g$ is a cyclic vector for $T$. After that, we characterize the frame property of such sequences in terms of the approximate point spectrum of $T^*$. Examples include certain operators (normal operators, compact operators, unilateral shifts, multiplication operators on Hardy spaces, etc.) that either generate only Riesz bases or allow redundancy. Our bridge theorem makes explicit the structural equivalence between frames generated by the functional calculus and function frames.
\end{abstract}


\keywords{2020 \emph{Mathematics subject classification.} 42C15; 47B02; 47A60; 47A10; 46C05.}
\keywords{\emph{Key words and phrases.} Hilbert space, frame, Riesz basis, operators, functional calculus, frames generated by the functional calculus, function frames, polynomial frames, dynamical frames, spectral theory, approximate point spectrum.}

\tableofcontents

\section{Introduction}

A central problem in frame theory is the construction and characterization of frames possessing additional algebraic or structural properties. Two particularly important classes naturally arise in this context:

\begin{itemize}
    \item[(i)] \textbf{Frames generated by the functional calculus}, that is, frames of the form
    \[
        \big(f(T)e_n\big)_{n\in\mathbb{N}},
    \]
    where $T\in \mathcal{B}(K)$, $f$ is a function defined on a neighborhood of the spectrum $\sigma(T)$, and $(e_n)_{n\in\mathbb{N}}$ is a sequence of vectors in $K$. 
    
    \item[(ii)] \textbf{Function frames}, that is, frames of the form
    \[
        \big(f_n(T)g\big)_{n\in\mathbb{N}},
    \]
    where $T\in \mathcal{B}(K)$, $(f_n)_{n\in\mathbb{N}}$ is a family of functions defined on a neighborhood of the spectrum $\sigma(T)$, and $g\in K$.
\end{itemize}

Function frames of the form $ \big(f_n(T)g\big)_{n\in\mathbb{N}}$ appear naturally in dynamical sampling theory, where one reconstructs signals from iterates of an evolution operator. This direction was developed in particular by Aldroubi–Cabrelli–Molter \cite{AldroubiCabrelliMolter2014,AldroubiCabrelliMolter2018}, Bownik–Speegle \cite{BownikSpeegle2015}, and later Christensen–Hasannasab \cite{ChristensenHasannasab2017-1, ChristensenHasannasab2017-2}. Frames generated via functional calculus belong to the broader theory of functional calculi, studied in \cite{Haase2006,Nikolski2002,Taylor1972}. 
The relationship between these two constructions, although implicit in the spectral theorem for normal operators \cite{Conway1990, NagyFoiasBercoviciKerchy2010, Nikolski2002}, does not seem to have been explicitly formulated in frame-theoretic terms.
The first contribution of this paper is to establish a rigorous and explicit bridge between these two fundamental constructions. \\ \\
\textbf{Bridge between frames generated by the functional calculus and function frames.} The first contribution of this paper can be summarized as follows: we establish a structural correspondence between frames generated by the functional calculus and function frames of a bounded normal operator (theorem~\ref{theorem-bridge}), showing that both constructions generate the ``same'' frames. This yields a conceptual unification that has not been explicitly formulated in the frame-theoretic literature.  \\ \\
\textbf{After that}, we derive a spectral characterization of such frames (theorem~\ref{thm:spectral}), formulated in terms of the approximate point spectrum of the underlying operator. This yields a criterion ensuring that sequences generated by the functional calculus are automatically Riesz bases whenever they form frames. From an operator-theoretic perspective, theorem \ref{thm:spectral} characterizes those operators for which surjectivity under holomorphic functional calculus forces invertibility. \\ \\
Together, these results contribute to the understanding of the analytic and spectral aspects underlying frames generated by the functional calculus and function frames and open new perspectives for applications in dynamical sampling, spectral theory, and functional analysis. \\ \\
\textbf{Plan of the paper.} Section \ref{SectionPreliminaries} establishes notation and recalls fundamental definitions and preliminary results on frames and spectral theory. Section \ref{SectionOrbitFrames} develops the theory of function frames and frames arising from the functional calculus, and establishes the connection between the two, when the operator is normal. Section \ref{MainSection} contains a spectral theorem on frames generated by the functional calculus and several detailed examples illustrating the theory.

\section{Preliminaries}
\label{SectionPreliminaries}

\subsection{Notation and conventions}
\label{SubsectionNotation}

Throughout this article, we adopt the following conventions:

$\mathbb{F}$ denotes either $\mathbb{R}$ or $\mathbb{C}$. Unless otherwise stated, all Hilbert spaces are complex and separable. 

A typical Hilbert space is denoted by $K$.

$\mathbb{N} = \{0,1,2,\ldots\}$ denotes the natural numbers including zero, while $\mathbb{N}^* = \mathbb{N} \setminus \{0\}$.

$\mathbb{D} = \{z \in \mathbb{C} : |z| < 1\}$ denotes the open unit disk, and $\overline{\mathbb{D}} = \{z \in \mathbb{C} : |z| \leq 1\}$ its closure.

For a subset $A$ of a topological space $X$, $\overline{A}$ denotes the closure, $\text{int}(A)$ the interior, and $\partial A$ the boundary.

For spectral sets:
\begin{itemize}
\item $\sigma(T)$ denotes the \textbf{spectrum} of $T$;
\item $\sigma_p(T) = \{\lambda \in \mathbb{C} : \Ker(T - \lambda I) \neq \{0\}\}$ denotes the \textbf{point spectrum} (eigenvalues);
\item $\sigma_{ap}(T) = \{\lambda \in \mathbb{C} : \inf_{\|x\|=1} \|(T-\lambda I)x\| = 0\}$ denotes the \textbf{approximate point spectrum};
\item $\rho(T) = \mathbb{C} \setminus \sigma(T)$ denotes the resolvent set.
\end{itemize}

For a measure space $(X,\Sigma,\mu)$, $L^2(X,\mu;\mathbb{F})$ denotes the Lebesgue space of square-integrable functions modulo equality $\mu$-almost everywhere.

For an operator $T \in \mathcal{B}(K)$, a vector $g \in K$ is called a cyclic vector for $T$ if $\overline{\Span \{T^n g\}_{n \in \mathbb{N}}} = K$.

\subsection{Frame theory}

Let $K$ be a separable Hilbert space.

\begin{definition}[Frames and related concepts]
\label{DefFrame}
A sequence $\Phi = (\varphi_n)_{n \in \mathbb{N}}$ in $K$ is called:
\begin{enumerate}[label=(\alph*)]
\item A \textbf{Bessel sequence} if there exists $B > 0$ such that
\[\forall v \in K : \sum_{n=0}^{\infty} |\langle v, \varphi_n \rangle|^2 \leq B \|v\|^2.\]
The infimum of all such $B$ is called the \textbf{Bessel bound}.

\item A \textbf{frame} if there exist constants $0 < A \leq B < \infty$ such that
\[\forall v \in K : A\|v\|^2 \leq \sum_{n=0}^{\infty} |\langle v, \varphi_n \rangle|^2 \leq B\|v\|^2.\]
The constants $A$ and $B$ are called \textbf{frame bounds}, and the optimal (largest $A$, smallest $B$) are the \textbf{optimal frame bounds}.

\item A \textbf{Riesz basis} if it is the image of an orthonormal basis under a bounded invertible operator, or equivalently, if $\Phi$ is a frame and $\{\varphi_n : n \in \mathbb{N}\}$ is linearly independent.

\item \textbf{Complete} if $\overline{\Span}\{\varphi_n : n \in \mathbb{N}\} = K$.
\end{enumerate}
\end{definition}

\begin{remark}
Every Riesz basis is a frame, but not every frame is a Riesz basis due to possible redundancy. A frame is a Riesz basis if and only if it is minimal, meaning that no element can be removed while preserving the frame property.
\end{remark}


\subsection{Spectral theory}

\begin{proposition}[Approximate point spectrum characterization]
\label{PropApproxSpecCharac}
For $T \in \mathcal{B}(K)$ and $\lambda \in \mathbb{C}$, the following are equivalent:
\begin{enumerate}[label=(\roman*)]
\item $\lambda \in \sigma_{ap}(T)$;
\item There exists a sequence $(x_n)_{n \in \mathbb{N}}$ in $K$ with $\|x_n\| = 1$ for all $n$ such that $\|(T - \lambda I)x_n\| \to 0$;
\item $T - \lambda I$ is not bounded below, i.e., $\inf_{\|x\|=1} \|(T-\lambda I)x\| = 0$.
\end{enumerate}
\end{proposition}

\begin{proposition}[Compactness of spectra]
$\sigma_{ap}(T)$ and $\sigma(T)$ are non-empty compact subsets of $\mathbb{C}$ when $K \neq \{0\}$.
\end{proposition}

\begin{proposition}[Spectral inclusions]
\label{PropSpectralInclusions}
For any $T \in \mathcal{B}(K)$:
\begin{enumerate}[label=(\roman*)]
\item $\sigma_p(T) \subseteq \sigma_{ap}(T) \subseteq \sigma(T)$;
\item $\sigma(T) = \sigma_{ap}(T) \cup \sigma_{ap}(T^*)^*$, where $A^* = \{\overline{z} : z \in A\}$;
\item $\partial \sigma(T) \subseteq \sigma_{ap}(T)$;
\item If $T$ is normal, then $\sigma_{ap}(T) = \sigma(T)$;
\end{enumerate}
\end{proposition}

\begin{theorem}[Holomorphic spectral mapping theorem]
\label{ThHolomorphicSpectralMapping}
Let $T \in \mathcal{B}(K)$ and let $f: \Omega \to \mathbb{C}$ be holomorphic on an open set $\Omega \supseteq \sigma(T)$. Then:
\[\sigma(f(T)) = f(\sigma(T)) = \{f(\lambda) : \lambda \in \sigma(T)\}.\]
\end{theorem}

\section{A connection between frames generated by the functional calculus and function frames in the case of a normal operator}
\label{SectionOrbitFrames}

\begin{remark}[Functional calculus]\label{rem:functional-calculus}
Throughout this article, we distinguish between three levels of functional
calculi:
\begin{itemize}
\item The \textbf{continuous functional calculus}, which applies to bounded
  normal operators and continuous functions on $\sigma(T)$, and suffices for
  the unitary equivalence results in the present section.
  
\item The \textbf{holomorphic functional calculus}, which is required in
  section \ref{MainSection}  to obtain spectral mapping properties for the approximate point
  spectrum and to use local factorization arguments of the form
  $f(z)-f(s)=(z-s)h(z)$.
  
  \item The \textbf{polynomial functional calculus}.
\end{itemize}
Unless otherwise stated, results involving spectral mapping properties of $f(T)$ are formulated under the assumption that $f$ is holomorphic
on an open neighborhood of $\sigma(T)$.
\end{remark}

\mbox{} \\
As mentioned in the introduction, there are two particularly important classes of frames that naturally arise. If we use the generic term ''function'' for a polynomial, holomorphic, or continuous function, then we have

\begin{itemize}
    \item[(i)] \textbf{Frames generated by the functional calculus}, that is, frames of the form
    \[
        \big(f(T)e_n\big)_{n\in\mathbb{N}},
    \]
    where $T\in \mathcal{B}(K)$, $f$ is a function defined on a neighborhood of the spectrum $\sigma(T)$, and $(e_n)_{n\in\mathbb{N}}$ is a sequence of vectors in $K$. 
    
    \item[(ii)] \textbf{Function frames}, that is, frames of the form
    \[
        \big(f_n(T)g\big)_{n\in\mathbb{N}},
    \]
    where $T\in \mathcal{B}(K)$, $(f_n)_{n\in\mathbb{N}}$ is a family of functions defined on a neighborhood of the spectrum $\sigma(T)$, and $g\in K$. \\
    This class comprises:
    \begin{itemize}
\item \textbf{Power orbit frames or dynamical frames \cite{AguileraCabrelliNegreiraPaternostro2025, AldroubiCabrelliMolter2014, AldroubiCabrelliMolter2018,BownikSpeegle2015,ChristensenHasannasab2017-1,ChristensenHasannasab2017-2, CabrelliMolterSuarez2024, MartinMedriMolter2021, Philipp2017}.} These are the frames of the type $(T^n g)_{n \in \mathbb{N}}$, obtained by considering $f_n(z) := z^n$ for all $n \in \mathbb{N}$. \\
It is known that for $(T^n g)_{n \in \mathbb{N}}$ to be a frame, $  S  $ should be a contraction ($  \Vert S \Vert \leq 1  $) with $  S^*  $ strongly stable, meaning $  (S^*)^n x \to 0  $ as $  n \to \infty  $ for all $  x \in K  $.
\item \textbf{Stricto sensu polynomial frames \cite{Dai2007, HasannasabKaldeweySchoppert2026,MhaskarNarcowichPrestinWard2000, Schoppert2025, Schoppert2026}}. These are true polynomials $(p_n)_{n \in \mathbb{N}}$ in $L^2(\mathbb{R}^n;\mathbb{R})$ (possibly restricted then to a subset of $\mathbb{R}^n$ such as the sphere $\mathbb{S}^{n-1}$ \cite{Dai2007,MhaskarNarcowichPrestinWard2000}) obtained by considering $f_n := p_n$ for all $n \in \mathbb{N}$, $T = Id$, and $g := 1$ the constant unit function. 
\end{itemize}
\end{itemize}
If these sequences are not frames, we simply call them \textbf{sequences generated by the functional calculus}, and \textbf{function sequences}, respectively. \\ \\
For bounded normal operators, we have the following spectral theorem

\begin{theorem}[Spectral theorem in cyclic form (see Conway (\cite{Conway1990}, Ch. X), Nikolski (\cite{Nikolski2002}, Ch. 1), or Sz.-Nagy–Foias \cite{NagyFoiasBercoviciKerchy2010})]
\label{theorem-spectral-decomposition-for-continuous-functional-calculus}
Let $K$ be a separable Hilbert space and $T \in B(K)$ a normal operator. 
Let $g \in K$ be a cyclic vector for $T$, i.e.
\[
\overline{\operatorname{span}}\{T^n g : n \ge 0\} = K.
\]
Then there exists a unique finite positive Borel measure $\mu_g$ on $\sigma(T)$ and a unitary operator
\[
U : L^2(\sigma(T), \mu_g) \to K
\]
such that
\[
U^* T U = M_z,
\]
where $(M_z \varphi)(z) = z \varphi(z)$.
Moreover,
\[
U(1) = g.
\]
\end{theorem}


\begin{proof}
We proceed in several steps. \\
\textbf{Step 1: Construction of the spectral measure $\mu_g$.} \\
Since $T$ is normal, the $C^*$-algebra $C^*(T,I)$ generated by $T$ and $I$ is commutative. 
By the continuous functional calculus for normal operators, there exists a *-isometric isomorphism
\[
\Phi : C(\sigma(T)) \longrightarrow C^*(T,I), 
\qquad f \mapsto f(T),
\]
satisfying $\Phi(z) = T$. \\
Define a linear functional $\Lambda_g$ on $C(\sigma(T))$ by
\[
\Lambda_g(f) := \langle f(T)g, g\rangle.
\]
Then \\
\begin{itemize}
\item $\Lambda_g$ is linear.
\item $\Lambda_g(f) \ge 0$ whenever $f \ge 0$, since $f(T)$ is positive.
\item $|\Lambda_g(f)| \le \|f(T)\|\|g\|^2 = \|f\|_\infty \|g\|^2$.
\end{itemize}
Thus $\Lambda_g$ is a positive bounded linear functional on $C(\sigma(T))$. \\
By the Riesz representation theorem, there exists a unique finite positive Borel measure $\mu_g$ on $\sigma(T)$ such that
\[
\langle f(T)g, g\rangle = \int_{\sigma(T)} f(z)\, d\mu_g(z)
\quad \text{for all } f \in C(\sigma(T)).
\]
\textbf{Step 2: Definition of the map $U$.} \\
Let $\mathcal{P}$ denote the polynomials. Define
\[
U_0 : \mathcal{P} \subseteq L^2(\sigma(T),\mu_g) \to K,
\qquad U_0(p) = p(T)g.
\]
For polynomials $p,q$, we compute
\[
\langle U_0(p), U_0(q)\rangle_K
= \langle p(T)g, q(T)g\rangle
= \langle (q\overline{p})(T)g, g\rangle
= \int_{\sigma(T)} p(z)\overline{q(z)}\, d\mu_g(z).
\]
Hence
\[
\langle U_0(p), U_0(q)\rangle_K
= \langle p, q\rangle_{L^2(\mu_g)},
\]
so $U_0$ is an isometry on polynomials. \\
Since polynomials are dense in $C(\sigma(T))$ and $C(\sigma(T))$ is dense in $L^2(\sigma(T),\mu_g)$, $U_0$ extends uniquely to an isometry
\[
U : L^2(\sigma(T),\mu_g) \to K.
\]
\textbf{Step 3: Surjectivity of $U$.} \\
The range of $U$ contains $\{p(T)g : p \text{ polynomial}\}$, whose closure equals $K$ by cyclicity of $g$. Since $U$ is an isometry, its range is closed. Hence $\operatorname{Ran}(U) = K$, so $U$ is unitary. \\
\textbf{Step 4: Intertwining relation.}
For a polynomial $p$,
\[
U^* T U p = U^* (T p(T)g) = U^*(z p(z))(T)g = z p(z).
\]
By density, this holds for all $\varphi \in L^2(\mu_g)$, so
\[
U^* T U = M_z.
\]
Finally, $U(1) = 1(T)g = g$.
\end{proof}


\begin{corollary}[Formula]
\label{corollary-image-under-U*-of-Tng-and-f(T)en}
Let $K$ be a separable Hilbert space, $T \in \mathcal{B}(K)$ be a normal operator, and $g \in K$ be a cyclic vector for $T$. Then $U^*(f(T)(e)) =  f \times U^* e$, and in particular $U^* g = 1$, for any function $f : \Omega \to \mathbb{C}$ with $\Omega \supseteq \sigma(T)$ and for any $e \in K$. \\
Consequently,
\begin{itemize}
\item For any sequence of continuous functions $(f_n)_{n \in \mathbb{N}}$ with $f_n : \Omega \to \mathbb{C}$ and $\Omega \supseteq \sigma(T)$, $(f_n(T)g)_{n \in \mathbb{N}}$ is a frame in $K$ if and only if $(f_n)_{n \in \mathbb{N}}$ is a frame in $L^2(\sigma(T), \mu_g)$.
\item For any continuous function $f: \Omega \to \mathbb{C}$ with $\Omega \supseteq \sigma(T)$ and any sequence of vectors $(e_n)$ of $K$, the sequence $(f(T)(e_n))$ is a frame in $K$ if and only if  $(f \times U^* e_n)$ is a frame in $L^2(\sigma(T),\mu_g)$.
\end{itemize}
\end{corollary}

\begin{proof}
\mbox{} \\
Recall from theorem \ref{theorem-spectral-decomposition-for-continuous-functional-calculus} that \(U^* T U = M_z \), so \(U^* f(T) U = f(M_z) = M_f\) (functional calculus intertwining). Therefore,
\[U^*(f(T)(e)) = U^* f(T) U (U^* e) = f(M_z)(U^* e) = f \times U^* e.\]
Since $U$ is unitary and \(U^* g = 1\) (deducible from the second step of the the proof of \ref{theorem-spectral-decomposition-for-continuous-functional-calculus}), the frame properties and bounds are preserved.
\end{proof}

The following bridge theorem can now be deduced.

\begin{theorem}[Sequences generated by the functional calculus as function sequences]
\label{theorem-bridge}
Let $K$ be a separable Hilbert space, $T \in \mathcal{B}(K)$ be a normal operator, and $g \in K$ be a cyclic vector for $T$. Then for any continuous function $f: \Omega \to \mathbb{C}$ with $\Omega \supseteq \sigma(T)$, we have
\begin{itemize}
\item $(f \times U^*e)(T)g = f(T)e$.
\item in particular: $f(T)g = Uf$.
\end{itemize}
Consequently
\begin{itemize}
\item  For any sequence of vectors $(e_n)_{n \in \mathbb{N}}$ in $K$ and for any continuous function $f: \Omega \to \mathbb{C}$ with $\Omega \supseteq \sigma(T)$, we have that $f(T)e_n = (f \times U^*e_n)(T)g$.  
\item For any sequence of continuous functions $(f_n)_{n \in \mathbb{N}}$ with $f_n : \Omega \to \mathbb{C}$ and $\Omega \supseteq \sigma(T)$, we have that $f_n(T)g = 1(T) Uf_n$.
\end{itemize}
\end{theorem}

\begin{proof}
For any function $f : \Omega \to \mathbb{C}$ with $\Omega \supseteq \sigma(T)$ and for any $e \in K$, we have, as a particular case of corollary \ref{corollary-image-under-U*-of-Tng-and-f(T)en}, that $U^*(f(T)g) = f \times U^* g = f$. Replacing $f \mapsto f \times U^* e$ gives $U^* \left[ (f \times U^*e) (T)g  \right]  = f \times U^* e$, and consequently $(f \times U^*e)(T)g = U(f \times U^*e) = f(T)e$, using corollary \ref{corollary-image-under-U*-of-Tng-and-f(T)en} again. For the second point, use corollary  \ref{corollary-image-under-U*-of-Tng-and-f(T)en} again with $e := g$: we have $U^*(f(T)g) = f$ and so $Uf = f(T)g$.
\end{proof}

\begin{remark}
\mbox{} \\
\begin{itemize}
\item Theorem \ref{theorem-bridge} means, in particular, that any frame generated by the functional calculus can be seen as a function frame, and vice versa. The spectral theorem is typically used to analyze operators, whereas here it is used to transfer structural properties of frames. This shift of viewpoint is the main conceptual contribution.

\item \textbf{On the role of cyclicity. } The assumption that $g$ is cyclic for $T$ is essential in theorems \ref{theorem-spectral-decomposition-for-continuous-functional-calculus} and \ref{theorem-bridge}.
Indeed, cyclicity ensures that the spectral representation induced by $g$
yields a unitary equivalence between $T$ and the multiplication operator $M_z$
on the whole space $L^2(\sigma(T),\mu_g)$. If $g$ is not cyclic, the map
\[
p \mapsto p(T)g
\]
defines an isometric embedding into a proper closed $T$-invariant subspace of
$\mathcal{K}$, and the associated representation only captures the restriction
of $T$ to the cyclic subspace generated by $g$.
In this case, the unitary equivalence holds only at the level of this cyclic
component, and the frame properties of $(f_n(T)g)$ must be analyzed via a
direct integral decomposition over the spectral multiplicity. \\
Thus, cyclicity is precisely the condition that allows one to reduce the study
of frames generated by $T$ to function frames on a scalar-valued $L^2$-space,
without loss of information.
\item In theorem \ref{theorem-bridge}, continuity of $f$ on $\sigma(T)$ is sufficient, since only the
continuous functional calculus is used. Unless additional conditions are assumed, the theorem fails to produce an equivalence between the two kinds of frames when the functions are polynomial or holomorphic and not merely continuous, since the function $U^* e_n$ is in general only continuous.
\end{itemize}
\end{remark}

\section{A spectral condition ensuring the frame property for sequences generated by the functional calculus}
\label{MainSection}

In this section, we establish that the condition $\sigma_{ap}(T^*) = \sigma(T^*)$ is necessary and sufficient for all sequences generated by the functional calculus of $T$ to be Riesz bases whenever they are frames. This result establishes that certain operators (normal operators, unilateral shifts, etc.) generate only Riesz bases, never redundant frames, while others allow redundancy.

\subsection{The criterion}

The following theorem provides a necessary and sufficient spectral condition that characterizes when frames generated by the functional calculus are automatically Riesz bases.

\begin{theorem}[Spectral characterization]
\label{thm:spectral}
Let $K$ be a separable Hilbert space and $T \in \mathcal{B}(K)$. Then the following are equivalent:
\begin{enumerate}[label=(\roman*)]
\item $\sigma_{ap}(T^*) = \sigma(T^*)$;
\item For every open set $\Omega \supseteq \sigma(T)$ and every holomorphic function $f: \Omega \to \mathbb{C}$, $f(T)$ is surjective implies $f(T)$ is invertible;
\item For every open set $\Omega \supseteq \sigma(T)$, every holomorphic function $f: \Omega \to \mathbb{C}$, and every orthonormal basis $(e_n)_{n \in \mathbb{N}}$ in $K$, $(f(T)(e_n))_{n \in \mathbb{N}}$ is a frame implies it is a Riesz basis.
\end{enumerate}
\end{theorem}

\begin{remark}
\mbox{} \\
\begin{itemize}
\item This theorem reveals a direct implication: the geometric property of frames being Riesz bases is completely determined by the spectral property $\sigma_{ap}(T^*) = \sigma(T^*)$. Moreover, this criterion is intrinsic to $T$ and independent of the choice of orthonormal basis or holomorphic function.
\item The relation between the approximate point spectrum and surjectivity is classical (see Conway \cite{Conway1990}, Ch. IV and Dunford–Schwartz \cite{DunfordSchwartz1958} Ch. V).
\end{itemize}
\end{remark}

\mbox{} \\
Before proving theorem \ref{thm:spectral}, we establish several key lemmas.

\begin{lemma}[Surjectivity and approximate point spectrum]
\label{Lemma(T-L)SurjApproxSpecAdjoint}
Let $K$ be a Hilbert space, $T \in \mathcal{B}(K)$, and $\lambda \in \mathbb{C}$. Then:
\[\overline{\lambda} \notin \sigma_{ap}(T^*) \Leftrightarrow T - \lambda I \text{ is surjective}.\]
\end{lemma}

\begin{proof}
($\Rightarrow$) Suppose $\overline{\lambda} \notin \sigma_{ap}(T^*)$. By definition, there exists $C > 0$ such that
\[\forall x \in K : \|T^*x - \overline{\lambda}x\| \geq C\|x\|.\]
This immediately implies $\Ker(T^* - \overline{\lambda}I) = \{0\}$. \\
Moreover, $\Ran(T^* - \overline{\lambda}I)$ is closed: if $y_n = (T^* - \overline{\lambda}I)(x_n)$ with $y_n \to y$, then
\[\|x_n - x_m\| \leq C^{-1}\|y_n - y_m\|,\]
so $(x_n)$ is Cauchy. Thus $x_n \to x$ for some $x \in K$, and by continuity, $y = (T^* - \overline{\lambda}I)(x)$. \\
By the closed range theorem, $\Ran(T - \lambda I)$ is also closed. Furthermore,
\[\overline{\Ran(T - \lambda I)} = \Ker(T^* - \overline{\lambda}I)^{\perp} = \{0\}^{\perp} = K.\]
Therefore, $T - \lambda I$ is surjective. \\
($\Leftarrow$) Suppose $T - \lambda I$ is surjective. Let $B$ be the Moore-Penrose pseudo-inverse of $T - \lambda I$, defined on $\Ran(T - \lambda I) = K$. Then $B(T - \lambda I) = P_{\Ker(T-\lambda I)^{\perp}}$, the orthogonal projection. \\
For any $x \in K$ with $\|x\| = 1$, write $x = x_1 + x_2$ where $x_1 \in \Ker(T - \lambda I)^{\perp}$ and $x_2 \in \Ker(T - \lambda I)$. Since $B$ exists and is bounded, there exists $\delta > 0$ such that
\[\|x_1\| \leq \delta\|(T - \lambda I)x_1\| = \delta\|(T - \lambda I)x\|.\]
Let $y \in K$ with $\|y\|=1$. By surjectivity of $T-\lambda I$, there exists $x \in K$ such that
\[
(T-\lambda I)x = y.
\]
Decompose $x = x_1 + x_2$ with $x_1 \in \Ker(T-\lambda I)^\perp$ and $x_2 \in \Ker(T-\lambda I)$. Then
\[
(T-\lambda I)x = (T-\lambda I)x_1,
\]
hence
\[
\|x_1\| \le \delta \|y\| = \delta.
\]
Now let $z \in K$. Using the identity
\[
\langle (T-\lambda I)x, z \rangle = \langle x, (T^*-\overline{\lambda}I)z \rangle,
\]
we obtain
\[
|\langle y, z \rangle|
= |\langle (T-\lambda I)x_1, z \rangle|
= |\langle x_1, (T^*-\overline{\lambda}I)z \rangle|
\le \|x_1\| \cdot \|(T^*-\overline{\lambda}I)z\|
\le \delta \|(T^*-\overline{\lambda}I)z\|.
\]
Taking the supremum over all $y$ with $\|y\|=1$, we deduce
\[
\|z\|
= \sup_{\|y\|=1} |\langle y,z\rangle|
\le \delta \|(T^*-\overline{\lambda}I)z\|.
\]
Thus
\[
\|(T^*-\overline{\lambda}I)z\| \ge \delta^{-1}\|z\|
\quad \text{for all } z \in K.
\]
This proves that $T^*-\overline{\lambda}I$ is bounded below, hence
\[
\overline{\lambda} \notin \sigma_{ap}(T^*).
\]
Thus $\overline{\lambda} \notin \sigma_{ap}(T^*)$. 
\end{proof}

\begin{remark}
The last argument relies on the closed range theorem and standard duality
properties of bounded operators; see, for instance, \cite[Chapter~IV]{Conway1990}
or \cite[Chapter~V]{DunfordSchwartz1958}
\end{remark}

\begin{lemma}[Spectral equivalence and surjectivity]
\label{LemmaApproxSpecEqualSpec}
Let $K$ be a Hilbert space and $T \in \mathcal{B}(K)$. Then:
\[\sigma_{ap}(T^*) = \sigma(T^*) \Leftrightarrow \left[\forall \lambda \in \mathbb{C} : T - \lambda I \text{ surjective} \Rightarrow T - \lambda I \text{ invertible}\right].\]
\end{lemma}

\begin{proof}
($\Rightarrow$) Assume $\sigma_{ap}(T^*) = \sigma(T^*)$ and let $\lambda \in \mathbb{C}$ be such that $T - \lambda I$ is surjective. By lemma \ref{Lemma(T-L)SurjApproxSpecAdjoint}, $\overline{\lambda} \notin \sigma_{ap}(T^*)$. Since $\sigma_{ap}(T^*) = \sigma(T^*)$, we have $\overline{\lambda} \notin \sigma(T^*)$. \\
Recall that $\sigma(T^*) = \overline{\sigma(T)}$ (the complex conjugate of $\sigma(T)$). Therefore, $\lambda \notin \sigma(T)$, which means $T - \lambda I$ is invertible. \\
($\Leftarrow$) Assume that for all $\lambda \in \mathbb{C}$, if $T - \lambda I$ is surjective, then $T - \lambda I$ is invertible. We prove $\sigma_{ap}(T^*) \supseteq \sigma(T^*)$ (the reverse inclusion always holds). \\
Let $\lambda \in \sigma(T^*) \setminus \sigma_{ap}(T^*)$. Then $\overline{\lambda} \in \sigma(T)$ and $\overline{\lambda} \notin \sigma_{ap}(T^*)$. By lemma \ref{Lemma(T-L)SurjApproxSpecAdjoint}, $T - \overline{\lambda}I$ is surjective. By hypothesis, $T - \overline{\lambda}I$ is invertible, so $\overline{\lambda} \notin \sigma(T)$, a contradiction.
\end{proof}

\begin{lemma}[Stability of spectral properties under holomorphic calculus]
\label{LemmaSpectrumFunctionOperator}
Let $K$ be a Hilbert space, $T \in \mathcal{B}(K)$, and $\Omega \supseteq \sigma(T)$ be open. Then:
\begin{enumerate}[label=(\roman*)]
\item If $\cl(\sigma_p(T)) = \sigma(T)$, then for every holomorphic $f: \Omega \to \mathbb{C}$, we have $\cl(\sigma_p(f(T))) = \sigma(f(T))$;
\item If $\sigma_{ap}(T) = \sigma(T)$, then for every holomorphic $f: \Omega \to \mathbb{C}$, we have $\sigma_{ap}(f(T)) = \sigma(f(T))$.
\end{enumerate}
\end{lemma}

\begin{proof}
(i) Assume $\cl(\sigma_p(T)) = \sigma(T)$ and let $f: \Omega \to \mathbb{C}$ be holomorphic. \\
Since $\sigma_p(f(T)) \subseteq \sigma(f(T))$ and $\sigma(f(T))$ is closed, we have $\cl(\sigma_p(f(T))) \subseteq \sigma(f(T))$. \\
Conversely, let $\lambda \in \sigma(f(T)) = f(\sigma(T))$ (by theorem \ref{ThHolomorphicSpectralMapping}). There exists $s \in \sigma(T)$ with $\lambda = f(s)$. \\
Since $\sigma(T) = \cl(\sigma_p(T))$, there exists a sequence $(s_n) \subseteq \sigma_p(T)$ with $s_n \to s$. For each $n$, choose $v_n \in K \setminus \{0\}$ with $Tv_n = s_n v_n$. \\ 
By the spectral mapping theorem for eigenvectors, $f(T)v_n = f(s_n)v_n$, so $f(s_n) \in \sigma_p(f(T))$. Since $f$ is continuous and $f(s_n) \to f(s) = \lambda$, we have $\lambda \in \cl(\sigma_p(f(T)))$. \\
(ii) Assume $\sigma_{ap}(T) = \sigma(T)$ and let $f: \Omega \to \mathbb{C}$ be holomorphic. \\
We always have $\sigma_{ap}(f(T)) \subseteq \sigma(f(T))$. For the reverse inclusion, let $\lambda \in \sigma(f(T)) = f(\sigma(T))$. Choose $s \in \sigma(T)$ with $\lambda = f(s)$. \\
Since $s \in \sigma_{ap}(T)$, there exists a sequence $(x_n)$ in $K$ with $\|x_n\| = 1$ and $(T - sI)x_n \to 0$. \\
Define $g: \Omega \to \mathbb{C}$ by $g(z) = f(z) - f(s)$. Since $g(s) = 0$ and $g$ is holomorphic, there exists a holomorphic function $h: \Omega \to \mathbb{C}$ such that $g(z) = h(z)(z - s)$ for all $z \in \Omega$ (by the removable singularity theorem in a neighborhood of $s$, extended to $\Omega$). \\
By the functional calculus, $g(T) = h(T)(T - sI)$. Therefore,
\[(f(T) - \lambda I)x_n = g(T)x_n = h(T)(T - sI)x_n \to 0\]
since $h(T)$ is bounded. Thus $\lambda \in \sigma_{ap}(f(T))$.
\end{proof}

\mbox{} \\
We also recall the following theorem.

\begin{theorem}[Frame sequences induced by operators {\cite[Theorem~5.5.4]{Christensen2016}}]
\label{thm:Christensen}
Let $\mathcal{K}$ be a Hilbert space, $A \in \mathcal{B}(\mathcal{K})$, and
$(e_n)_{n\in\mathbb{N}}$ an orthonormal basis of $\mathcal{K}$. Then:
\begin{enumerate}
  \item $(Ae_n)_{n\in\mathbb{N}}$ is a frame if and only if $A$ is surjective.
  \item $(Ae_n)_{n\in\mathbb{N}}$ is a Riesz basis if and only if $A$ is invertible.
\end{enumerate}
\end{theorem}

\mbox{} \\
We now prove our main theorem of this section.

\begin{proof}[Proof of theorem \ref{thm:spectral}]
We prove (i) $\Leftrightarrow$ (ii) $\Leftrightarrow$ (iii). \\
\textbf{(i) $\Rightarrow$ (ii):} Assume $\sigma_{ap}(T^*) = \sigma(T^*)$. Let $f: \Omega \to \mathbb{C}$ be holomorphic with $\Omega \supseteq \sigma(T)$. \\
By lemma \ref{LemmaSpectrumFunctionOperator}(ii) applied to $T^*$, we have $\sigma_{ap}(f(T)^*) = \sigma(f(T)^*)$. By lemma \ref{LemmaApproxSpecEqualSpec} applied to $f(T)$, if $f(T)$ is surjective, then $f(T)$ is invertible. \\
\textbf{(ii) $\Rightarrow$ (iii):} Assume (ii) holds. Let $(e_n)_{n \in \mathbb{N}}$ be an orthonormal basis and suppose $(f(T)(e_n))_{n \in \mathbb{N}}$ is a frame. By theorem \ref{thm:Christensen}, $f(T)$ is surjective. By (ii), $f(T)$ is invertible. Again by theorem \ref{thm:Christensen}, $(f(T)(e_n))_{n \in \mathbb{N}}$ is a Riesz basis. \\
\textbf{(iii) $\Rightarrow$ (i):} Assume (iii) holds. For any $\lambda \in \mathbb{C}$, consider $f(z) = z - \lambda$ defined on $\Omega = \mathbb{C} \supseteq \sigma(T)$. Then $f(T) = T - \lambda I$. Let $(e_n)$ be any orthonormal basis. \\
Suppose that $T - \lambda I$ is surjective. By theorem \ref{thm:Christensen}, $(f(T)(e_n))_{n \in \mathbb{N}} = ((T - \lambda I)(e_n))_{n \in \mathbb{N}}$ is a frame. By (iii), it is a Riesz basis, so $T - \lambda I$ is invertible. \\
By lemma \ref{LemmaApproxSpecEqualSpec}, $\sigma_{ap}(T^*) = \sigma(T^*)$.
\end{proof}

\subsection{Applications and examples}
\label{SectionApplications}

The class of operators satisfying $\sigma_{ap}(T^*) = \sigma(T^*)$ includes normal operators \cite{Halmos1961} (like bilateral shifts), compact operators, unilateral shifts, and multiplication operators on Hardy spaces.

\begin{example}[Normal operators]
\label{ExBilateralShift}
\mbox{} \\
Let $K = \ell^2(\mathbb{Z})$ with orthonormal basis $(e_n)_{n \in \mathbb{Z}}$. The bilateral shift $T: \ell^2(\mathbb{Z}) \to \ell^2(\mathbb{Z})$ is defined by $T(e_n) = e_{n+1}$. Then 
\begin{enumerate}[label=(\roman*)]
\item $T$ is unitary;
\item $\sigma(T) = \sigma_{ap}(T) =\partial \mathbb{D}$ (the unit circle);
\item $\sigma_p(T) = \emptyset$ (no point spectrum);
\item $\sigma_{ap}(T^*) = \sigma(T^*) = \partial \mathbb{D}$;
\end{enumerate}
Now, by theorem \ref{thm:spectral}, for any holomorphic $f: \Omega \to \mathbb{C}$ with $\Omega \supseteq \overline{\mathbb{D}}$, $(f(T)(e_n))_{n \in \mathbb{Z}}$ is a frame if and only if it is a Riesz basis if and only if $f$ has no zeros on $\partial \mathbb{D}$, 
\end{example}

\begin{example}[Compact operators]
\mbox{} \\
Consider $K = \ell^2(\mathbb{N})$ and define $T$ by $T(e_0) = 0$ and $T(e_n) = \frac{1}{n} e_{n-1}$ for $n \geq 1$. Then
\begin{enumerate}[label=(\roman*)]
\item $T$ is compact since it is a Hilbert-Schmidt operator;
\item $\sigma(T) = \sigma_{ap}(T) = \sigma_p(T) = \{0\}$;
\item $\sigma_{ap}(T^*) = \sigma(T^*) = \{0\}$;
\item Since $f(T)$ is compact and $\ell^2(\mathbb{N})$ is infinite-dimensional, $f(T)$ is not surjective.
\end{enumerate}
This example illustrates that compact operators in infinite dimensions always satisfy our criterion but are never surjective, so they provide frames only in the simple finite-dimensional cases.
\end{example}

\begin{example}[Unilateral shifts \cite{Halmos1961,Nikolski2002}]
\label{ExWeightedShift}
\mbox{} \\
Let $K = \ell^2(\mathbb{N})$ with canonical basis $(e_n)_{n \in \mathbb{N}}$. Define the right shift $S: \ell^2(\mathbb{N}) \to \ell^2(\mathbb{N})$ by
\[S(e_n) =  e_{n+1}, \quad n \in \mathbb{N}.\]
Then:
\begin{enumerate}[label=(\roman*)]
\item $\sigma(S) = \overline{\mathbb{D}}$; 
\item $\sigma_{ap}(S) = \partial \mathbb{D}$;
\item $\sigma_p(S) = \emptyset$.
\item $\sigma_{ap}(S^*) = \sigma(S^*) = \overline{\mathbb{D}}$.
\end{enumerate}
By theorem \ref{thm:spectral}, for the unilateral shift $S$
\begin{itemize}
\item $(e_n + e_{n+1})_{n \in \mathbb{N}}$ corresponds to $f(z) = 1 + z$, which has $f(-1) = 0 \in f(\overline{\mathbb{D}})$, so it is not a frame;
\item More generally, $(e_n + e_{n+1} + \cdots + e_{n+k})_{n \in \mathbb{N}}$ corresponds to $f(z) = \frac{1 - z^{k+1}}{1 - z}$, which has roots in $\overline{\mathbb{D}}$ (the $(k+1)$-th roots of unity), so it is not a frame;
\item $(ae_n + be_{n+1})_{n \in \mathbb{N}}$ with $a, b \in \mathbb{C}$ corresponds to $f(z) = a + bz$, which is a frame iff $|a| > |b|$ (so that $0 \notin f(\overline{\mathbb{D}})$), in which case it is a Riesz basis.
\end{itemize}
\end{example}

\begin{example}[Multiplication operators on Hardy spaces \cite{GarciaMashreghiRoss2016}]
\label{ExHardySpace}
Let $H^2(\mathbb{D})$ denote the Hardy space of holomorphic functions on $\mathbb{D}$ with
\[\|f\|^2 = \sup_{0 < r < 1} \frac{1}{2\pi} \int_0^{2\pi} |f(re^{i\theta})|^2 \, d\theta < \infty.\]
The multiplication operator $M_z: H^2(\mathbb{D}) \to H^2(\mathbb{D})$ defined by $(M_z f)(z) = zf(z)$ satisfies:
\begin{enumerate}[label=(\roman*)]
\item $\sigma(M_z) = \overline{\mathbb{D}}$;
\item $\sigma_{ap}(M_z^*) = \sigma(M_z^*) = \overline{\mathbb{D}}$;
\end{enumerate}
So, by theorem \ref{thm:spectral}, for any holomorphic $f: \Omega \to \mathbb{C}$ with $\Omega \supseteq \overline{\mathbb{D}}$ and orthonormal basis $(e_n)_{n \in \mathbb{N}}$ of $H^2(\mathbb{D})$, $(f(T)(e_n))_{n \in \mathbb{Z}}$ is a frame if and only if it is a Riesz basis if and only if $f$ has no zeros on $\partial \mathbb{D}$. \\
Since the monomials $(z^n)_{n \in \mathbb{N}}$ form an orthonormal basis of $H^2(\mathbb{D})$, we can conclude that $(f(M_z)(z^n))_{n \in \mathbb{N}}$ is a frame if and only if it is a Riesz basis if and only if $f$ has no zeros in $\overline{\mathbb{D}}$. \\
This connects our results to classical complex analysis and operator theory on Hardy spaces.
\end{example}


%\section*{Acknowledgments}
%
%The author thanks the anonymous referees for their careful reading and valuable suggestions that improved the presentation of this work.

\nocite{*}
\bibliographystyle{plain}
\bibliography{references}

\Addresses

\end{document}
